\chapter{Shapes and Space}\label{ch:g1:shapes}

Windows are rectangles, wheels are circles, roofs hide triangles:
shapes are everywhere. This chapter learns to recognize them by
counting sides and corners, and to say precisely \emph{where} things
are: left, right, above, below, between.

\section{The four basic shapes}

\begin{definition}[Square, rectangle, triangle, circle]\label{def:g1:shapes:def}
\begin{itemize}
  \item the \emph{triangle} has $3$ sides and $3$ corners;
  \item the \emph{square} has $4$ equal sides and $4$ corners;
  \item the \emph{rectangle} has $4$ sides and $4$ corners, with
        opposite sides equal (a stretched square);
  \item the \emph{circle} is perfectly round: no sides, no corners.
\end{itemize}
\end{definition}

\begin{omfigure}
\begin{tikzpicture}[scale=0.7]
  \draw[thick, fill=omDef!12] (0,0) -- (2,0) -- (1,1.7) -- cycle;
  \node[below, align=center, font=\small] at (1,-0.25)
    {triangle:\\$3$ sides};
  \begin{scope}[xshift=3.6cm]
  \draw[thick, fill=omThm!12] (0,0) rectangle (1.7,1.7);
  \node[below, align=center, font=\small] at (0.85,-0.25)
    {square:\\$4$ equal sides};
  \end{scope}
  \begin{scope}[xshift=6.9cm]
  \draw[thick, fill=omMeth!12] (0,0.05) rectangle (2.8,1.65);
  \node[below, align=center, font=\small] at (1.4,-0.25)
    {rectangle:\\$4$ sides};
  \end{scope}
  \begin{scope}[xshift=11.3cm]
  \draw[thick, fill=omProp!12] (0.9,0.85) circle (0.9);
  \node[below, align=center, font=\small] at (0.9,-0.25)
    {circle:\\no corner};
  \end{scope}
\end{tikzpicture}

{\small The four basic shapes. To name a shape, do not trust the
first glance: \emph{count} its sides and corners.}
\end{omfigure}

\begin{example}[Tilted shapes are still shapes]\label{ex:g1:shapes:tilted}
A square standing on its corner is still a square --- four equal
sides, four corners --- even if it looks like a kite. Turning a
shape does not change what it is.
\end{example}

\begin{example}[Shape hunt]\label{ex:g1:shapes:hunt}
In the classroom: the board is a rectangle, the clock a circle, a
folded napkin can be a triangle, some windows are squares. Real
objects are not perfect shapes, but their faces are close to them
(solids and their faces return in \cref{ch:g2:shapes}).
\end{example}

\section{Where is it?}

\begin{definition}[Position words]\label{def:g1:shapes:position}
To say where something is: \emph{left} / \emph{right},
\emph{above} / \emph{below}, \emph{on} / \emph{under},
\emph{in front of} / \emph{behind}, \emph{between},
\emph{inside} / \emph{outside}.
\end{definition}

\begin{omfigure}
\begin{tikzpicture}[scale=0.75]
  \draw[thick, fill=omDef!12] (0,0) circle (0.55);
  \draw[thick, fill=omThm!12] (1.9,-0.55) rectangle (3.0,0.55);
  \draw[thick, fill=omMeth!12] (4.4,-0.5) -- (5.5,-0.5) -- (4.95,0.45)
    -- cycle;
  \node[below, font=\small] at (0,-0.75) {circle};
  \node[below, font=\small] at (2.45,-0.75) {square};
  \node[below, font=\small] at (4.95,-0.75) {triangle};
  \node[right, font=\small, align=left] at (6.3,0)
    {the square is \emph{between}\\the circle and the triangle;\\
    the circle is \emph{left} of the square};
\end{tikzpicture}

{\small Position words in action. Careful: \emph{your} left and the
left of someone facing you are opposite!}
\end{omfigure}

\section{Paths on a grid}

\begin{method}[Describing a path]\label{met:g1:shapes:path}
On grid paper, a path is a list of steps:
$\to$ (one square right), $\leftarrow$ (left), $\uparrow$ (up),
$\downarrow$ (down).
\begin{enumerate}
  \item To follow a path: start on the marked square and obey the
        arrows one by one;
  \item to describe a path: walk it in your head and write the
        arrows in order;
  \item two different paths can join the same two squares --- some
        shorter, some longer.
\end{enumerate}
\end{method}

\begin{omfigure}
\begin{tikzpicture}[scale=0.6]
  \draw[gray!40] (0,0) grid (7,4);
  \fill[omDef!40] (0.5,0.5) +(-0.42,-0.42) rectangle +(0.42,0.42);
  \node[font=\footnotesize] at (0.5,0.5) {start};
  \fill[omThm!40] (5.5,2.5) +(-0.42,-0.42) rectangle +(0.42,0.42);
  \node[font=\footnotesize] at (5.5,2.5) {goal};
  % seven unit arrows, one per step, with small gaps
  \foreach \a/\b in {0.5/1.5, 1.5/2.5}
    \draw[-Stealth, omDef, very thick] (\a+0.1,0.5) -- (\b-0.1,0.5);
  \foreach \a/\b in {0.5/1.5, 1.5/2.5}
    \draw[-Stealth, omDef, very thick] (2.5,\a+0.1) -- (2.5,\b-0.1);
  \foreach \a/\b in {2.5/3.5, 3.5/4.5}
    \draw[-Stealth, omDef, very thick] (\a+0.1,2.5) -- (\b-0.1,2.5);
  \draw[-Stealth, omDef, very thick] (4.6,2.5) -- (5.05,2.5);
\end{tikzpicture}

{\small A path from start to goal, one arrow per step:
$\to\ \to\ \uparrow\ \uparrow\ \to\ \to\ \to$ --- seven steps. Can
you find another path with the same number of steps?}
\end{omfigure}

\section{Exercises}

\begin{exercise}[$\star$]\label{exo:g1:shapes:1}
For each shape, count the sides and corners and name it: a shape
with $3$ corners; a round shape with no corner; a shape with $4$
equal sides.
\end{exercise}

\begin{exercise}[$\star$]\label{exo:g1:shapes:2}
Find at home: two rectangles, one circle, one triangle, one square.
(Doors? plates? road signs?)
\end{exercise}

\begin{exercise}[$\star$]\label{exo:g1:shapes:3}
True or false? ``A square turned on its corner stops being a
square.'' Explain with \cref{ex:g1:shapes:tilted}.
\end{exercise}

\begin{exercise}[$\star$]\label{exo:g1:shapes:4}
Draw on grid paper: a square of $3$ squares by $3$; a rectangle of
$5$ by $2$; a triangle using a diagonal of the grid.
\end{exercise}

\begin{exercise}[$\star$]\label{exo:g1:shapes:5}
Place three toys in a row. Describe the row with position words:
which is between the others? Which is on the left?
\end{exercise}

\begin{exercise}[$\star$]\label{exo:g1:shapes:6}
Draw a circle. Draw a cross \emph{inside} it, a star
\emph{outside} it, and a dot \emph{on} the circle itself.
\end{exercise}

\begin{exercise}[$\star$]\label{exo:g1:shapes:7}
On grid paper, mark a start square. Follow:
$\to\ \uparrow\ \to\ \uparrow\ \to$. Mark the arrival. Then write
the path that goes straight back to the start.
\end{exercise}

\begin{exercise}[$\star$]\label{exo:g1:shapes:8}
How many corners in total: two triangles and one square? One circle
and two rectangles?
\end{exercise}

\begin{exercise}[$\star\star$]\label{exo:g1:shapes:9}
Take $12$ matches. Can you outline a square using \emph{all} of
them, with whole matches and the same number on each side? How many
matches per side? Could you do it with $10$ matches?
\end{exercise}
